\section{Registros de acarreo e incidencia y reconstrucción proyectiva de historias}
\label{lib07:armonizacion}
La construcción APP--TRIT--TPK fija un estado enriquecido antes de
sus realizaciones numéricas, incidenciales y operatorias. Los registros
que intervienen en esa construcción se distinguen por la operación
que permiten invertir y por las relaciones que conservan. Esta sección
reúne esas distinciones y compone la recuperación con el refinamiento.

\subsection{Acarreo, historia, incidencia y archivo}
El acarreo de las secciones~\ref{lib03:residuo} y
\ref{lib03:memoria} es un dato aritmético entero. Su conservación,
junto con la profundidad y el estado final, restituye el antecedente
por la inversa explícita~\eqref{lib03:inversa-afin}.
La historia de transporte contiene la sucesión ordenada de operaciones,
hojas, orientaciones y retornos. En ella, dos etapas con igual fase
pueden tener registros acumulados distintos. La holonomía actúa sobre
esa historia y sobre el marco en que se efectúan sus lecturas.

El registro de incidencia conserva las relaciones de frontera que
utilizan las construcciones excepcionales. El objeto ordenado \(K\)
y su lectura \(\kappa_{\rm ag}\) intervienen con las normalizaciones
y los mapas de la sección~\ref{lib01:identificacion}.
La suficiencia para recuperar un operador se expresa allí mediante
las respuestas del marco y la inversión correspondiente.
La memoria de archivo de la sección~\ref{lib04:archivo}, por su parte,
es una sucesión de componentes de un espacio de Hilbert. La identidad
telescópica~\eqref{lib04:telescopia} conserva la norma entre el estado
remanente y las componentes archivadas; la inversa
\eqref{lib04:recuperacion-infinita} reconstruye la amplitud inicial.

Estas cuatro funciones se componen conservando sus tipos.
Un entero de acarreo, una palabra de transporte, una incidencia
y una amplitud archivada constituyen datos distintos del sistema.
Las pruebas de recuperación especifican qué combinación de ellos
es suficiente para cada lectura. La figura~\ref{lib07:fig-registros}
representa esa jerarquía funcional.

\begin{figure}[htbp]
\centering
\begin{tikzpicture}[
 node distance=9mm,
 box/.style={draw=ink,rounded corners=2pt,fill=white,align=center,
 text width=58mm,minimum height=15mm,font=\small},
 arr/.style={-{Latex[length=2mm]},draw=ink,thick}]
\node[box,text width=125mm] (s) {Estado enriquecido APP--TRIT--TPK\\
 operación, hoja, orientación, frontera y memoria};
\node[box] (a) at (-33mm,-30mm)
 {Acarreo aritmético\\Residuo, profundidad e inversa};
\node[box] (h) at (33mm,-30mm)
 {Historia de transporte\\Composición ordenada y holonomía};
\node[box] (q) at (-33mm,-62mm)
 {Registro de incidencia\\Lecturas de frontera y marco};
\node[box] (m) at (33mm,-62mm)
 {Archivo de recuperación\\Componentes y reconstrucción};
\node[box,text width=125mm] (r) at (0,-94mm)
 {Recuperación compatible con el refinamiento\\
 Inversas tipadas y relaciones conservadas};
\draw[arr] (s.south -| a.north)--(a);
\draw[arr] (s.south -| h.north)--(h);
\draw[arr] (a)--(q);
\draw[arr] (h)--(m);
\draw[arr] (q.south)--(r.north -| q.south);
\draw[arr] (m.south)--(r.north -| m.south);
\end{tikzpicture}
\caption{Organización de los registros según su función.
Las flechas inferiores indican composición de condiciones de
recuperación, y no identificación de los cuatro tipos de dato.}
\label{lib07:fig-registros}
\end{figure}

\subsection{Árbol de refinamiento y pesos locales}
Sea \(H_k\) el conjunto finito de prefijos admisibles de profundidad
\(k\), con truncamientos \(p_{k+1,k}\), del sistema de historias ya
construido. Una emisión se cuenta con su multiplicidad y
\(d(\zeta)>0\) denota el número de emisiones desde \(\zeta\).
La distribución uniforme de las semillas y el peso local
\(1/d(\zeta)\) definen \(\mu_k\). La suma de los pesos de las
prolongaciones de un prefijo coincide con su peso: las medidas son
compatibles y determinan la medida cilíndrica \(\mu\).
Se demuestran a continuación las proposiciones de memoria de fibra y
naturalidad de la medida, con estas definiciones explícitas.

\subsection{Lectores enriquecidos y memoria de fibra}
\label{sec:ch-memoria-fibra}

La diferencia entre historia y valor admite un criterio operativo. En un
nivel finito \(H_k\), sea \(q_k:H_k\to Z_k\) un lector y sea
\(M_k:H_k\to\mathcal M_k\) el registro conservado. El lector enriquecido
es \(\widehat q_k=(q_k,M_k)\). La memoria distingue estados de una misma
fibra del lector; su suficiencia se comprueba en el dominio considerado.
Esta formulación incorpora al propio artículo el desarrollo completo de la
información de fibra que se utiliza a continuación.

\begin{proposition}[Recuperabilidad en un nivel finito]
\label{prop:ch-recuperabilidad-fibra}
Existe una inversa de \(\widehat q_k\) sobre su imagen si y sólo si
\(\widehat q_k\) es inyectivo. Para una variable aleatoria \(Z\) con
valores en \(H_k\) y distribución \(\mu_k\), se cumple
\[
 \mathsf H(Z)=\mathsf H(q_k(Z))+\mathsf H(Z\mid q_k(Z)).
\]
Si \(\widehat q_k\) es inyectivo sobre el soporte de \(Z\), entonces
\[
 \mathsf H(Z\mid q_k(Z),M_k(Z))=0.
\]
Aquí \(\mathsf H(Z)=-\sum_z\Pr(Z=z)\log\Pr(Z=z)\), con
\(0\log0=0\), es la entropía finita en nats; la entropía condicional es
el promedio de la entropía de las distribuciones condicionadas.
\end{proposition}

\begin{proof}
Una inversa recupera un único antecedente, lo que exige inyectividad;
recíprocamente, ésta permite asignar a cada par publicado su único
antecedente. Puesto que \(q_k(Z)\) es función determinista de \(Z\),
agrupar la suma que define \(\mathsf H(Z)\) por fibras de \(q_k\) da la
regla de la cadena. Si el par \((q_k(Z),M_k(Z))\) distingue el soporte,
cada distribución condicionada correspondiente está concentrada en un solo
estado y su entropía es cero.
\end{proof}

No publicar una coordenada y borrar el registro que permite recuperarla son
operaciones distintas. Este resultado no identifica entropía con cardinalidad
del continuo ni introduce una interpretación térmica: su dominio, lector y
distribución son los que se acaban de fijar.

\subsection{Equivarianza de la medida bajo simetrías del árbol}
\label{sec:ch-naturalidad-medida}

\begin{proposition}[Transporte conjunto del árbol y de su medida]
\label{prop:ch-naturalidad-medida}
Sea \(g=(g_k)_{k\geq0}\) una familia de biyecciones \(g_k:H_k\to H_k\)
que conmuta con los truncamientos, fija la raíz y transporta biyectivamente
las emisiones admisibles, conservando sus multiplicidades. Entonces
\[
 d(g_k\zeta)=d(\zeta),\qquad
 (g_k)_*\mu_k=\mu_k,\qquad (g_\infty)_*\mu=\mu.
\]
Más generalmente, si \(\mu_\zeta\) es la ley de continuaciones de una
historia \(\zeta\in H_k\), se tiene
\[
 (g_\infty)_*\mu_\zeta=\mu_{g_k\zeta}.
\]
\end{proposition}

\begin{proof}
La biyección de emisiones conserva su número con multiplicidades. La
distribución uniforme de semillas también es invariante. Cada trayectoria
y su imagen reciben, por tanto, el mismo producto de pesos locales
\(1/d(\zeta)\). Se obtiene la igualdad de medidas en cada nivel y en
cada cilindro. La coherencia con los truncamientos define \(g_\infty\),
y la unicidad de la medida determinada por cilindros prueba la igualdad
en el límite. El mismo argumento, comenzando en \(\zeta\), prueba la
afirmación sobre continuaciones.
\end{proof}

Se obtiene así la naturalidad del árbol y de la medida: un cambio coherente
de representación transporta los pesos y no obliga a escogerlos de nuevo.
La conservación de emisiones y multiplicidades es una hipótesis esencial de
la proposición, no un dato añadido después de la construcción.

\subsection{Recuperación de historias y transporte de pertenencia}
\begin{proposition}[Compatibilidad de inversas y familias]
\label{lib07:prop-familias}
Sean \(H_k\) los niveles anteriores y \(Y_k=\widehat q_k(H_k)\).
Supóngase que \(\widehat q_k:H_k\to Y_k\) es biyectivo y que existen
truncamientos \(s_{k+1,k}\) tales que
\[
 s_{k+1,k}\widehat q_{k+1}=\widehat q_k p_{k+1,k}.
\]
Entonces las inversas finitas inducen una biyección
\[
 \widehat q_\infty:\varprojlim H_k\longrightarrow\varprojlim Y_k.
\]
Si \(H=\varprojlim H_k\), \(Y=\varprojlim Y_k\) y
\(\mathcal T(A)=\widehat q_\infty[A]\), la aplicación
\(\mathcal T:\mathcal P(H)\to\mathcal P(Y)\) es biyectiva y
preserva uniones, intersecciones, complementos relativos y pertenencia.
Cada restricción \(A\to\mathcal T(A)\) es una biyección.
\end{proposition}
\begin{proof}
Al componer la identidad de compatibilidad con las inversas se obtiene
\[
 p_{k+1,k}\widehat q_{k+1}^{-1}
   =\widehat q_k^{-1}s_{k+1,k}.
\]
Por tanto, una familia compatible \(y=(y_k)\) determina la familia
compatible \(x=(\widehat q_k^{-1}(y_k))\). Las dos aplicaciones
coordenadas son inversas. La imagen directa por una biyección tiene
como inversa la imagen directa por la biyección inversa. Para toda
familia \((A_i)_{i\in I}\) de subconjuntos de \(H\),
\[
 \begin{aligned}
 \mathcal T\Bigl(\bigcup_i A_i\Bigr)&=\bigcup_i\mathcal T(A_i),\\
 \mathcal T\Bigl(\bigcap_i A_i\Bigr)&=\bigcap_i\mathcal T(A_i),\\
 \mathcal T(H\setminus A)&=Y\setminus\mathcal T(A).
 \end{aligned}
\]
La primera igualdad usa la definición de imagen. Para la segunda y
la tercera, la unicidad del antecedente permite trasladar las
condiciones de pertenencia. Finalmente,
\(x\in A\) equivale a
\(\widehat q_\infty(x)\in\mathcal T(A)\), lo que prueba también
la afirmación sobre restricciones.
\end{proof}

El teorema de refinamiento de rutas~\ref{lib05:teo-rutas}
proporciona una realización de esta compatibilidad sobre las etiquetas
aritméticas y sus acarreos. El resultado distingue la recuperación
de cada estado y la conservación de las operaciones sobre familias.
En el nivel de amplitudes, el levantamiento unitario
de~\ref{lib05:prop-unitaria} transporta además el producto interno.
En el nivel probabilístico, la proposición
\ref{prop:ch-naturalidad-medida} conserva los pesos cuando se
transportan las emisiones con sus multiplicidades.

\subsection{Articulación con la clausura genealógica del continuo}
La sección~\ref{sec:cierre-cardinal-nativo} construye una clausura
genealógica \(\mathfrak G\) y establece la identidad cardinal
\(2^{\aleph_0}=\aleph_1\) sobre su potencia interna. La clausura,
el código generador y la historia realizada forman parte de su
definición; su representación conjuntista se establece después de la
construcción genealógica, mediante las operaciones y demostraciones
incorporadas en dicha sección.

La presente articulación incorpora las pruebas de memoria,
naturalidad y transporte de familias que utiliza esa relación con
la conservación. La cadena demostrativa queda discriminada:
las inversas recuperan estados; su compatibilidad recupera
historias; la biyección transporta las operaciones de pertenencia;
la clausura genealógica de la sección~\ref{sec:cierre-cardinal-nativo} fija el dominio
del enunciado cardinal. Así, la recuperación y la organización
de familias se integran en la doble proyección con sus respectivos
objetos, operaciones y conclusiones.
